\documentclass[11pt]{article}
\usepackage[margin=1in]{geometry}
\usepackage{amsmath,amssymb}

\title{Disproof of Conjecture 00000001061\\
\large Equal discriminants and different point-count parities}
\author{}
\date{}

\begin{document}
\maketitle

Conjecture 00000001061 claims that the quadratic-character sign of the
discriminant determines the parity of a hyperelliptic point count. Two
smooth genus-two curves over the same field give opposite parities while
having equal discriminants.

\section*{The two curves}

Work over \(k=\mathbb F_5\). For \(c\in\{0,-1\}\), set
\[
  f_c(x)=x^5-x+c,\qquad C_c:\ y^2=f_c(x),
\]
where \(C_c\) denotes the smooth projective model. The derivative
\(f'_c=-1\) shows that the quintic is squarefree over an algebraic closure
of \(k\). Its simple roots also show that it is a nonsquare in
\(\overline{k}(x)\). Thus the double cover defined by \(y^2=f_c(x)\)
is geometrically integral.

The second chart, with coordinates \(u=1/x\) and \(v=y/x^3\), has equation
\[
  v^2=u-u^5+c u^6.
\]
Both charts are finite over their respective affine charts of
\(\mathbb P^1\), since the equations are monic in \(y\) and \(v\).
Their gluing is finite over \(\mathbb P^1\) and hence projective.
The derivative of \(y^2-f_c(x)\) with respect to \(x\) is \(1\),
so the first chart is smooth. In the second chart, the only point
over \(u=0\) is \((u,v)=(0,0)\), where the derivative of
\(v^2-u+u^5-cu^6\) with respect to \(u\) is \(-1\), so that point is
smooth as well. This also proves there is exactly one \(k\)-point at
infinity on each curve.

For the genus, the degree-two map to \(\mathbb P^1\) branches precisely
over the five simple zeros of \(f_c\) and at infinity. Each has
ramification index two, since the respective right-hand side has a
simple zero in its chart. At every other point the derivative in the
square-root variable is nonzero. The ramification is tame in
characteristic five, so Riemann--Hurwitz gives
\(2g(C_c)-2=2(-2)+6=2\), and \(g(C_c)=2\).

\section*{Discriminants and point counts}

The resultant of a monic degree-five polynomial with the constant
polynomial \(-1\) is the product of five copies of \(-1\), one for each
root counted with multiplicity. Hence both quintics have discriminant
\[
  \operatorname{disc}(f_c)=(-1)^{5\cdot4/2}
     \operatorname{Res}(f_c,f'_c)=(-1)^5=4 \quad\text{in }k.
\]
The quadratic character of \(4\) is \(+1\), since \(4=2^2\)
in \(k\). The discriminant polynomial in \(y\) of \(y^2-f_c(x)\) is
\(4f_c(x)\), whose polynomial discriminant is also the same for the
two curves, by the degree-five scaling factor \(4^8\).

For every \(x\in k\), \(x^5=x\), so \(f_c(x)=c\). Above each of the
five choices of \(x\), \(C_0\) has the unique affine point with \(y=0\).
Above each choice of \(x\), \(C_{-1}\) has two affine points, with
\(y=2\) and \(y=-2\). Including the respective points at infinity yields
\[
  \#C_0(k)=5+1=6,\qquad \#C_{-1}(k)=10+1=11.
\]
Thus the same field, genus, and discriminant give different parities.
The affine counts \(5\) and \(10\) also have different parities. Euler's
criterion correctly evaluates the common character sign, but that sign
cannot determine either parity.

\section*{Lean verification}

The accompanying project uses Lean 4.19.0 and the standard library.
It computes the formal derivative from the coefficient list, evaluates the
constant-polynomial resultant, and enumerates the five field elements in
each affine fibre and at infinity. The two final theorems rule out a
parity predictor from the discriminant sign on this family, for projective
and affine counts respectively. Run \verb|lake build| in \verb|lean/| to
verify both theorems and print their axiom dependencies.

\end{document}
